% 子文件：B_几何_初中
% 本文件包含 8 个 section


%% ==============================
%% 第9题：尺规作图确定点 F 并求解内错角
%% ==============================
\newpage

\section{解答：作 AD 垂直平分线与 $\angle BAC$ 角平分线}

\vspace{0.6cm}

\begin{center}
\begin{tikzpicture}[>=Stealth, line join=round, line cap=round, scale=1.2]
  % 1. 定义基本参数
  \def\a{3.2} % 正方形边长

  % 2. 定义定点坐标
  \coordinate (A) at (0, \a);
  \coordinate (B) at (0, 0);
  \coordinate (C) at (\a, 0);
  \coordinate (D) at (\a, \a);
  \coordinate (E) at (\a/2, \a);
  
  % 3. 绘制已知原始几何图形 (正方形及对角线)
  \draw[line width=1.0pt] (A) -- (B) -- (C) -- (D) -- cycle;
  \draw[line width=1.0pt] (A) -- (C);
  
  % 标注正方形顶点及 E 点
  \node[left, font=\large] at (A) {$A$};
  \node[left, font=\large] at (B) {$B$};
  \node[right, font=\large] at (C) {$C$};
  \node[right, font=\large] at (D) {$D$};
  \node[above right, font=\large, xshift=-2pt] at (E) {$E$};
  
  % ----------------------------------------------------
  % 第一步：尺规作图作 AD 的垂直平分线 l
  % ----------------------------------------------------
  \def\Rone{2.2} % 所用圆的半径 (要求 > 0.5*a)
  \pgfmathsetmacro{\dy}{sqrt(\Rone*\Rone - (\a/2)*(\a/2))}
  \coordinate (P1) at (\a/2, \a + \dy); % 上方相交点
  \coordinate (P2) at (\a/2, \a - \dy); % 下方相交点
  
  % 利用 clip 裁剪显示手工作图的交叉弧线痕迹
  \begin{scope}
    \clip (P1) circle (0.4);
    \draw[line width=0.6pt] (A) circle (\Rone);
    \draw[line width=0.6pt] (D) circle (\Rone);
  \end{scope}
  \begin{scope}
    \clip (P2) circle (0.4);
    \draw[line width=0.6pt] (A) circle (\Rone);
    \draw[line width=0.6pt] (D) circle (\Rone);
  \end{scope}
  
  % 绘制贯穿的垂直平分线 l
  \draw[line width=1.2pt] (\a/2, \a+1.6) -- (\a/2, -1.0) node[left, near start, font=\large] {$l$};
  
  % ----------------------------------------------------
  % 第二步：尺规作图作角 BAC 的平分线，并确定交点 F
  % ----------------------------------------------------
  \def\Rtwo{1.5} % 第一道弧半径
  % 在 \angle BAC 范围内画弧
  \draw[line width=0.6pt] (A) ++(-90:\Rtwo) arc (-90:-45:\Rtwo);
  
  % 获取角两边上的交点
  \coordinate (K1) at (0, \a-\Rtwo);
  \coordinate (K2) at ({\Rtwo*cos(-45)}, {\a+\Rtwo*sin(-45)});
  
  % 从 K1, K2 出发画第二道弧，形成交叉点 X
  \def\Rthree{2.0} % 相交弧半径
  % 计算大致交点以进行视口剪裁。角平分线方向为 -67.5 度
  % 距 A 点的理论距离 = \Rtwo * cos(22.5) + sqrt(\Rthree^2 - (\Rtwo*sin(22.5))^2)
  \pgfmathsetmacro{\Dbg}{\Rtwo*cos(22.5) + sqrt(\Rthree*\Rthree - (\Rtwo*sin(22.5))*(\Rtwo*sin(22.5)))}
  \coordinate (X) at ($(A) + (-67.5:\Dbg)$);
  
  \begin{scope}
    \clip (X) circle (0.5);
    \draw[line width=0.6pt] (K1) circle (\Rthree);
    \draw[line width=0.6pt] (K2) circle (\Rthree);
  \end{scope}
  
  % 计算平分线与垂直平分线 l 的理论交点 F
  % F的x坐标为 \a/2。其离A的水平距离为 \a/2。斜率为 tan(-67.5) = - (1 + sqrt(2)) = -2.41421356
  \pgfmathsetmacro{\yF}{\a - (\a/2)*2.41421356}
  \coordinate (F) at (\a/2, \yF);
  
  % 绘制角平分线 (从 A 穿过 X 和 F)
  % 延长一些以符合题图的连线特征
  \draw[line width=1.0pt] (A) -- ($(A)!1.08!(F)$);
  \node[right, font=\large] at (F) {$F$};

\end{tikzpicture}
\end{center}

\vspace{0.6cm}

{\small\color{gray}
\textbf{解题说明：}\\
1. \textbf{作图分析（中垂线与角平分线）：}\\
(1) 题目要求“作 $AD$ 的垂直平分线 $l$”，利用尺规作图：分别以定点 $A, D$ 为圆心，以大于 $\frac{1}{2}AD$ 为半径画弧，两弧在上方和下方各交于两点，用直尺且连接这两个交点，得到垂直平分线 $l$。显然由于 $ABCD$ 是正方形，$AB$ 垂直于 $AD$，因此垂直平分线 $l$ 必平行于 $AB$。\\
(2) 题目要求“点 $F$ 到 $\angle BAC$ 两边的距离相等”，由初中几何的\textbf{角平分线逆定理}可知，如果一个点到角两边的距离相等，则该点一定在该角的角平分线上。所以这实际上是要求我们作出 $\angle BAC$ 的角平分线。画弧痕迹：以 $A$ 为圆心任意长画弧交两边，再以两个交点为圆心同半径画弧，两小段弧相交形成小叉号。连接 $A$ 与该叉号交点并向外延伸射线便得到角平分线，其与中垂线 $l$ 碰撞相交得到的点即为要求的最终悬空点 $F$。\\
2. \textbf{第（2）小题的角度计算过程：}\\
已知四边形 $ABCD$ 是正方形，则其每个内角为 $90^\circ$（即 $\angle BAD = 90^\circ$）。\\
且 $AC$ 为正方形内部的长对角线，所以平分了原直角 $\angle BAC = \frac{1}{2}\angle BAD = 45^\circ$。\\
通过刚才的尺规作图得知直线 $AF$ 为 $\angle BAC$ 的精确角平分线，因此分出的内部夹小角 $\angle BAF = \frac{1}{2}\angle BAC = 22.5^\circ$。\\
又因为直线 $l$ 是外边线段 $AD$ 的垂直平分线，且方框内部侧边本身就满足 $AB \perp AD$。由于同一平面内垂直于同一条直线的两直线必平行可强力推导出：直线 $l \parallel AB$（即含有目标点的射线段 $EF \parallel AB$）。\\
最后，依据平行线的基本性质：“两条直线平行，被第三条直线所截，形成的内错角相互等同”。我们凝视由这两根绝对平行纵线夹出的“Z”字形翻折内角段，故直接稳健得出最终所求目标值： $\angle EFA = \angle BAF = 22.5^\circ$。
}

\vspace{0.6cm}

{\small\color{blue!70!black}
\textbf{绘图基本原理与步骤：}\\
\textbf{1. 几何底层参数化渲染：}先确立具有边长 \texttt{\textbackslash def\textbackslash a\{3.2\}} 参数的正方形主框架基座，用 \texttt{coordinate} 方法锁定全部核心顶点以杜绝硬编码死写数字，完全保证后续尺规线矢量拓扑稳定缩放结构。\\
\textbf{2. 严格的代数视窗裁剪逼真复原圆轨痕迹：}摒弃手动勾勒不可靠自由短弧的低级画法。对于中垂线的长跨度关键相交弧，利用宏代数 \texttt{\textbackslash pgfmathsetmacro} 结合勾股算式算出绝对理论中心交点位置距，动态将其自身设定为全包围小锚点，外裹原生态命令 \texttt{\textbackslash clip} 开启圆形蒙版隔离，随后在其背部罩满完整的画圆圆环 \texttt{circle} 坐标指令。利用计算机渲染切除特性极度自然逼真地只给画面剔透流出“手划针脚相抵交印痕”的残留图线。在 $\angle BAC$ 分水岭的角平分两段交锋处同样调用 \texttt{clip} 方法完成高精度原核模拟同理复刻。\\
\textbf{3. 精确极线方程推导规避连线穿模错配：}结尾处，需通过作图线把实点 $A$ 连接飞出指派至悬空末端交接点 $F$ 时；通过前述物理几何论证其必锁定位于 $x=a/2$ 的绝对中垂系统轨道上；叠加角平分线原始自带方向导数对应的绝对几何极偏转角为 $-67.5^\circ$，求解析正切极限值 $\tan(-67.5^\circ) = -(1+\sqrt{2}) \approx -2.4142$。由此极坐标切距值回送给代数机 \texttt{pgfmathsetmacro} 便可极度精准地测前向后下发预测该交接端点的高度差与截面绝对悬空标高坐标。完美严丝合缝接轨中垂轨段线，物理封杀凭感觉盲连导向极易突发的乱位溢出透视错穿等假象。
}

%% ==============================
%% 第7题：方格纸中的图形旋转、面积计算与格点确定
%% ==============================
\newpage

%% ==============================
%% 第7题：方格纸中的图形旋转、面积计算与格点确定
%% ==============================
\newpage



%% ==============================
%% 第15题：最短路径问题（饮马问题及模型应用）
%% ==============================
\newpage

\section{解答：最短路径问题（饮马模型及应用）}

\vspace{0.6cm}

\textbf{（1）如图①，在直线 $\boldsymbol{l}$ 上找一点 $\boldsymbol{P}$，使得 $\boldsymbol{PA+PB}$ 最短（画图即可）。}

\begin{center}
\begin{tikzpicture}[>=Stealth, line join=round, line cap=round, scale=1.2]
  % 1. 基本结构：直线 l 和两个点 A, B
  \draw[line width=1.0pt, black] (0, 0) -- (6, 0) node[right, font=\large] {$l$};
  \coordinate (A) at (1.5, 1.0);
  \coordinate (B) at (4.5, 2.5);
  
  \fill (A) circle (2pt) node[above left] {$A$};
  \fill (B) circle (2pt) node[above right] {$B$};
  
  % 2. 作图步骤：作点 A 关于直线 l 的对称点 A'
  \coordinate (A1) at (1.5, -1.0); % A 的对称点
  \fill[blue!70!black] (A1) circle (2pt) node[below right, text=blue!70!black] {$A'$}; % 将 A' 画出并加注蓝色区分作图印记
  \draw[dashed, blue!70!black, line width=0.8pt] (A) -- (A1); % 对称辅助虚线
  
  % 在横轴投影视角画一组直角标记
  \draw[blue!70!black, line width=0.6pt] (1.5, 0.2) -- (1.7, 0.2) -- (1.7, 0);
  
  % 3. 找点 P：连接 A' 和 B，交直线 l 于点 P
  % 运用 TikZ 的智能交点功能 intersect 找到 l 和射线的交汇
  \coordinate (P) at (intersection of A1--B and 0,0--6,0);
  \fill[red!80!black] (P) circle (2.5pt) node[below right, text=red!80!black, yshift=-3pt] {$P$};
  
  % 绘制辅助连接线 A'B
  \draw[dashed, blue!70!black, line width=0.8pt] (A1) -- (B);
  
  % 4. 绘制原像上的最终极值最短折线段路线： AP, PB
  \draw[red!80!black, line width=1.5pt] (A) -- (P) -- (B);
  
  % 图①编号
  \node[below, font=\large] at (3, -1.8) {①};
\end{tikzpicture}
\end{center}

\vspace{0.4cm}

\textbf{（2）如图②，在正方形 $\boldsymbol{ABCD}$ 中，$\boldsymbol{E}$ 为 $\boldsymbol{AB}$ 的中点，在线段 $\boldsymbol{BD}$ 上找一点 $\boldsymbol{P}$，使得 $\boldsymbol{PA+PE}$ 的值最小（画图即可）。}

\begin{center}
\begin{tikzpicture}[>=Stealth, line join=round, line cap=round, scale=1.3]
  \def\a{3.0} % 正方形基本变长
  % 1. 基本结构：原正方形及给定的中心对角连线图元
  \coordinate (A) at (0, \a);
  \coordinate (D) at (\a, \a);
  \coordinate (B) at (0, 0);
  \coordinate (C) at (\a, 0);
  
  \draw[line width=1.0pt, black] (A) -- (D) -- (C) -- (B) -- cycle;
  
  % 绘制原题已经含有的对角线段 AC 和 BD
  \draw[line width=0.8pt, black] (A) -- (C);
  \draw[line width=0.8pt, black] (B) -- (D);
  
  \node[above left] at (A) {$A$};
  \node[above right] at (D) {$D$};
  \node[below left] at (B) {$B$};
  \node[below right] at (C) {$C$};
  
  % 定义出左侧一半位置上的点 E
  \coordinate (E) at (0, \a/2);
  \fill (E) circle (2pt) node[left] {$E$};
  
  % 2. 饮马模型原理活用化归：
  % 点 A 取法相对要求反弹的最底折面（轴是对角线 BD）自然得出的绝对系统对称点，正是其自身位于全完美方框底围实相存在的另外一条右下角端点 C！
  
  % 利用算法找到连接 EC 这段线与既底大长轴 BD 的横切物理碰撞点即为我们要求的全局总长最短中转 P
  \coordinate (P) at (intersection of E--C and B--D);
  \fill[red!80!black] (P) circle (2.5pt) node[above, text=red!80!black, yshift=2pt] {$P$};
  
  % 绘制映射替代的最省长辅助视线 CE
  \draw[dashed, blue!70!black, line width=1.0pt] (E) -- (C);
  
  % 3. 压高渲染原始模型题干实际所走的半段路线 (A -> P -> E)
  \draw[red!80!black, line width=1.5pt] (A) -- (P) -- (E);
  
  % 图②编号
  \node[below, font=\large] at (\a/2, -0.6) {②};
\end{tikzpicture}
\end{center}


\vspace{0.4cm}

\textbf{（3）如图③，在正方形 $\boldsymbol{ABCD}$ 中，$\boldsymbol{E}$ 为 $\boldsymbol{AB}$ 的中点，动点 $\boldsymbol{M}$ 在 $\boldsymbol{BC}$ 上，动点 $\boldsymbol{N}$ 在 $\boldsymbol{CD}$ 上，连接 $\boldsymbol{EM, MN, NA}$，请你应用 (1) 的原理，在图③中找出点 $\boldsymbol{M, N}$，使得 $\boldsymbol{EM+MN+NA}$ 的值最小（画图即可）。}

\begin{center}
\begin{tikzpicture}[>=Stealth, line join=round, line cap=round, scale=1.3]
  \def\a{3.0} % 正方形主骨架宽度设定
  % 1. 全局基本黑视大边框构建
  \coordinate (A) at (0, \a);
  \coordinate (D) at (\a, \a);
  \coordinate (B) at (0, 0);
  \coordinate (C) at (\a, 0);
  
  \draw[line width=1.0pt, black] (A) -- (D) -- (C) -- (B) -- cycle;
  
  \node[above left] at (A) {$A$};
  \node[above right] at (D) {$D$};
  \node[below left] at (B) {$B$};
  \node[below right] at (C) {$C$};
  
  \coordinate (E) at (0, \a/2);
  \fill (E) circle (2pt) node[left] {$E$};
  
  % 2. “双管齐下”两段式延展镜像作法逻辑：
  % 分别作源路径起始头 E 和 末段收缩端 A 关于所在反射折返边壁（即底面底线 BC 、与对面强身墙线 CD）做各自的双向对外镜像发散点投影。
  
  % 第一维展开作案：建立 E 相对横亘底轴 BC 的下翻镜像 E' (- \a/2)
  \coordinate (E1) at (0, -\a/2);
  \fill[blue!70!black] (E1) circle (2pt) node[left, text=blue!70!black] {$E'$};
  \draw[dashed, blue!70!black, line width=0.8pt] (E) -- (E1); % 辅助垂直投影拉扯虚线段
  
  % 在横轴投影视角画一组直角标记
  \draw[blue!70!black, line width=0.6pt] (0.2, 0) -- (0.2, -0.2) -- (0, -0.2);
  
  % 第二维发力拓界：确立原始收口 A 点相对于另一强反射立面 CD 沿出边界横插到平面的绝对平移反位点 A'
  \coordinate (A1) at (2*\a, \a);
  \fill[blue!70!black] (A1) circle (2pt) node[above, text=blue!70!black] {$A'$};
  \draw[dashed, blue!70!black, line width=0.8pt] (A) -- (A1); % 横向飞拉线
  
  % 在纵轴投影视角画一组直角标记
  \draw[blue!70!black, line width=0.6pt] (\a, \a-0.2) -- (\a+0.2, \a-0.2) -- (\a+0.2, \a);
  
  % 3. 终结大势已定：使用两仪终极解——暴力直连已被我们破开完全平展化为一条平直极通线 E'A'
  \draw[dashed, blue!70!black, line width=1.0pt] (E1) -- (A1);
  
  % 4. 寻觅目标点：利用其虚直线划破实心原始二维底区时撞破割断两堵物理边界的原始出入角落差定位出真实题解 M 和 N。
  \coordinate (M) at (intersection of E1--A1 and B--C);
  \coordinate (N) at (intersection of E1--A1 and C--D);
  
  \fill[red!80!black] (M) circle (2.5pt) node[above left, text=red!80!black, xshift=-1pt, yshift=2pt] {$M$};
  \fill[red!80!black] (N) circle (2.5pt) node[below left, text=red!80!black, xshift=-2pt, yshift=-4pt] {$N$};
  
  % 5. 大功告成将原主命题极具弯曲折叠性质且无解般的内部路线真实浮现高亮填色
  \draw[red!80!black, line width=1.5pt] (E) -- (M) -- (N) -- (A);
  
  % 图③编号
  \node[below, font=\large] at (\a/2, -0.6-\a/2) {③};
\end{tikzpicture}
\end{center}

\vspace{0.6cm}
\hrule
\vspace{0.4cm}

{\small\color{gray}
\textbf{解题解析说明：}\\
这三道小题同属于经典的初中数学同构压轴模型：\textbf{“将军饮马”的最短路径问题及其多阶衍生。}其核心思想就在于利用点对线段（镜面）的对称原理将折线段拉扯并展开映射到一整侧，最后化归为最基本的两点间直线位移。\\
1. \textbf{第（1）题（基础一次反射）：}题目要求在给定平直边界 $l$ 上找到打在岸边的拐点 $P$ 使得入与出路径加和 $PA+PB$ 最小。由“两点之间线段最短”，只需找出 $A$ 相对边界 $l$ 的垂直镜像对称点 $A'$，然后用直尺直接连接 $A'B$ 两点。由于垂直平分线性质约束了该线段必与 $l$ 交过一个点且长度就是全域最极致的小，从而找出交点定为 $P$ 即可。\\
2. \textbf{第（2）题（正方形巧妙应用）：}要求在固定的斜坡对角线 $BD$ 上寻走点 $P$ 使得两线夹带 $PA+PE$ 最短。将此置于有着极强高轴对称度正方形 $ABCD$ 框架中。无需像上题那般盲目标划辅助线建立：点 $A$ 关于约束转折轴对角线 $BD$ 因为正方形自身的属性自然生成的绝对物理镜像点——正是图内的另一固定实景定点 $C$。以此题立马被化解成了单纯在问求“$PC+PE$”的最短折线值。由于定点均已暴露无遗，把固定起点的 $C$ 与定死身位的 $E$ 大笔连接，其一刀斩截 $BD$ 直线所遗留的必然是极正解 $P$。\\
3. \textbf{第（3）题（高难度二次复合反弹推图）：}本题动点横跨两条且彼此完全垂直的硬性反射面（下端硬底 $BC$ 面与右侧边沿墙角线 $CD$），构建这套诸如台球桌角连续反弹两遭所累积的最少全损耗：此时只需发挥 (1) 原理的三生万物的威力，让最开头的初始端发力点 $E$ 对脚下那堵墙 $BC$ 实施狠狠一次下翻强发展开成完全镜像出外援点 $E'$；然后同理针对落脚点的终极大头目的地 $A$ ，继续毫不手软地对着隔壁那一堵垂直的墙壁（触板 $CD$）如法炮制硬向反域展平一格平移出一个对等的假定目的地 $A'$。所有点皆开界后，所有障碍墙尽破，两极连点相通成毫无弯曲阻挠且必是最短路径的一整条长达虚拟走廊大直线（$E'A'$）。顺藤摸瓜回头看其在经过虚拟位面前物理戳爆了主平原两次护栏边界割裂下来的那两具血淋淋的节点交界便是正解题里呼之欲出的全息打点 $M$ 和 $N$ 所在。
}

\vspace{0.6cm}

{\small\color{blue!70!black}
\textbf{绘图基本原理与步骤：}\\
\textbf{1. 锚点基础代数建模：}在建立这三组同气连枝的小图结构中其基本的定位 \texttt{coordinate} 分散坐标系统均摒弃生硬硬算，把整个全系定边基元用极强普适扩展性的参数 \texttt{\textbackslash a} 进行依赖耦合。比如第三幅的对称目标只需要推写下 \texttt{A1(2*\textbackslash a, \textbackslash a)} 即可轻松跨过所有比例运算约束完成精准大框架对标对立。\\
\textbf{2. 代数交点原研求解器（\texttt{intersection of}）：}全代码组的核心黑科技解压神器！使用 TikZ 系统自带且在此物理线相交题中能够发挥核弹级效率的自带计算命令。在精细且直观地推演算出假想镜像辅助端点的相对偏移外延参数后（即使像第三张跨过全底沿这种极限折平外抛投射），不用手工推解各种斜率联立解析二元甚至二次多元方程根。直观交给机器运行 \texttt{coordinate (M) at (intersection of E1--A1 and B--C);} 让图形库算法直接在虚拟内存截下空间完美投射切割心。保证每次成图必然重合并且具备小数点满位精确度。\\
\textbf{3. 分层风格与虚实图纸分立术：}针对此题考研强烈的空间物理拉扯脑暴：在维持原版基础的墨黑 $1.0\rm{pt}$ 基调中；利用浅发冷调又伴有散漫感的蓝色 \texttt{dashed} 点划色系将抽象推理出所有平移投射出来的多维虚拟折射镜像区、辅外投引线作完全退位后方的大背书参照展示，既保留完整的思维解答推演论证走线全细节却又一点不抢夺本源戏份：最后仅给题纸提问的最重点答案大纲（所行最优质有效折算路轨）以最为极尽粗大的正红高警预报属性描出 \texttt{red!80!black} 及 \texttt{line width=1.5pt} 并追加强加注了全部解题落位黑点！
}

%% ==============================
%% 第16题：爱心捐赠活动条形与扇形统计图
%% ==============================
\newpage



%% ==============================
%% 第15题：以 O 为顶点的中心对称直角三角形的绘制
%% ==============================
\newpage

\section{解答：作对称直角三角形}

\vspace{0.6cm}

\textbf{15. 如图，已知点 $O$，以 $O$ 为直角顶点，作一个直角三角形以及该直角三角形的对称图形。}

\vspace{0.4cm}

\begin{center}
\begin{tikzpicture}[>=Stealth, line join=round, line cap=round]
  
  % 定义直角三角形的直角边长 \w=3.0, \h=2.0 
  % 并且保持图纸空间留白，因此适当用 \scale 可以放大缩小
  \def\w{2.5}
  \def\h{2.0}

  % 原点 O
  \coordinate (O) at (0, 0);

  % ===============================
  % 原直角三角形 (第一象限方向)
  % ===============================
  % 直角固定在 O 点
  \coordinate (A1) at (\w, 0);
  \coordinate (B1) at (0, \h);
  % 绘制红色的原图形线条
  \draw[red!80!black, line width=1.2pt] (O) -- (A1) -- (B1) -- cycle;

  % ===============================
  % 关于 O 点的中心对称图形 (第三象限方向)
  % ===============================
  \coordinate (A2) at (-\w, 0);
  \coordinate (B2) at (0, -\h);
  % 绘制红色的对称镜像线条
  \draw[red!80!black, line width=1.2pt] (O) -- (A2) -- (B2) -- cycle;

  % ===============================
  % 维持原题中原本的点与字标识
  % ===============================
  % 原本黑色的实体源点
  \fill[black] (O) circle (2.0pt);
  
  % 以带透明灰度感渲染大写代号 O，使得字母稍微避让重红色线圈 
  \node[font=\large, gray!80!black, below left, xshift=-1pt, yshift=-1pt] at (O) {$O$};

\end{tikzpicture}
\end{center}

\vspace{0.6cm}
\hrule
\vspace{0.4cm}

{\small\color{gray}
\textbf{解题分析与说明：}\\
1. \textbf{解析题意：}题目要求“以 $O$ 为直角顶点作一个直角三角形”，意味着这个三角形必须包含 $90^\circ$ 的直角并且将其顶点死锁在原试卷给出的 $O$ 黑色实心点上。随后题干后半句“以及该直角三角形的\textbf{对称图形}”虽然未指明对称轴或是对称中心的名字，但根据仅给出了一个唯一几何条件原点 $O$ 的语境分析可知，这里考察的恰恰必须是：\textbf{直角三角形围绕 $O$ 点发生 $180^\circ$ 全方位旋转倒装变换后的中心对称图形}（同时也是沿角分线的重轴对称折射图形）。\\
2. \textbf{绘图路径：}我们在定点 O 的右侧平直方向及正上方各取两截正交的垂置线段（例如底角长 $2.5$，竖高 $2.0$ 相对比例），形成一个纯向着“第一象限”生长的右上方直角三角形底模；接下来，利用反向定心等长延长策略，将它的实边向反方向分别长驱直入地探穿过 $O$ 点跑到正左方与正下方，再次将其双反端点相连接，最后能产生一个完全朝向“第三象限”倒挂生长的绝对底置镜像三角形，这就绘制出了极其经典的领结型中心全对称结构。
}

\vspace{0.6cm}

{\small\color{blue!70!black}
\textbf{绘图（制图）基本原理与步骤：}\\
\textbf{1. 极简代数相对位移建模：}本题为了呈现最为标准、不存在一点人工像素毛刺的严谨工程几何答卷，完全放弃在图纸上网格尺规硬抠点位。开篇立即利用极其干净自带泛型编程属性的两个自适应矢量宏边长参数 \texttt{\textbackslash def\textbackslash w\{2.5\}} 与 \texttt{\textbackslash def\textbackslash h\{2.0\}} 将所要张开的那张无形大网的长宽比例完全固死。构筑第一基础直角时，直接向空间锚射出 \texttt{(\textbackslash w, 0)} 及 \texttt{(0, \textbackslash h)} 构筑原像面；而在构筑反转全等对称副本时，系统层面上仅仅非常粗暴干脆地将绝对坐标轴前置强加负号正负突变为 \texttt{(-\textbackslash w, 0)} 及 \texttt{(0, -\textbackslash h)} ，保证了在非人眼干涉的纯数学运算推导层面上，不偏不倚地完成了毫无偏移的极致精确 $180^\circ$ 心锚全倒翻转反相。\\
\textbf{2. 隔离叠涂透胶渲染术：}这道制图的两枚大三角形主体轮廓采用统一粗长连贯而绝无破绽断桥的封闭式线组 \texttt{-{}-cycle} 闭合强拉算法，通过一支强力笔触直接描出了标准答案那高对比度亮眼的红色预警边沿 \texttt{draw[red!80!black, line width=1.2pt]}；而在作画完毕的最顶层即由于多条粗硬锐利红线条挤压在一起最为逼仄恶劣的原核心焦点处，最后加装覆盖回一颗带原生质感的原生小深黑心圈，以及将其斜下角的那个原名牌 $O$ 调成透着灰色薄雾带灰度的 `gray` 且施加微小向量阻尼侧滑 \texttt{xshift=-1pt} 的位移动画。以力挽狂澜维持住最初题干里最初那种冷淡清越孤点的原始试卷底稿原始物理风貌属性，不由于作答的大张旗鼓而被破坏磨灭。
}

%% ==============================
%% 第6题：不等号解集的数轴展现
%% ==============================
\newpage


%% ==============================
%% 第30题（补充题）：尺规作图（角平分线与垂直平分线）
%% ==============================
\newpage

\section{解答：尺规作图综合}

\vspace{0.6cm}

\textbf{\LARGE 15.} \textbf{如图，用直尺和圆规完成下列作图（保留作图痕迹）：}

\vspace{0.4cm}

\textbf{解：}

\textbf{（1）作图原理与底层执行逻辑：}

\textbf{原理剥析：} 已知原 $\triangle ABC$ 中 $\angle C = 90^\circ$且 $\angle B = 60^\circ$。要求在底边 $AC$ 上建立点 $D$，使得局部 $\triangle BCD$ 的三个内角凑齐 $30^\circ$、$60^\circ$、$90^\circ$。因为 $\triangle BCD$ 已经不可避免地占用了 $\angle C = 90^\circ$，那么剩下的另外两个内角必须强制锁定为 $30^\circ$ 和 $60^\circ$。由于原生母角 $\angle ABC$ 恰好为 $60^\circ$，因此寻找满足条件的 $D$ 点，本质上就是\textbf{要求我们严格去作 $\angle ABC$ 的角平分线}。一旦将其平分开来，顶部的 $\angle CBD$ 即被精准劈成 $30^\circ$，此时根据内角和法则，下方 $D$ 节点产生的 $\angle CDB$ 自然就被锁定为了 $180^\circ - 90^\circ - 30^\circ = 60^\circ$。
\textbf{尺规步进实录：}
\begin{enumerate}
    \item 以顶点 $B$ 为圆心，将圆规张开一定长度画下第一道测试弧，将原线段 $BC$ 与 $BA$ 分别斩断，留下两个刻骨交点。
    \item 随后拔出圆规，以上述两个交点为次级爆发圆心，以大于这二者真实间距一半的稳固长度为新半径画追踪弧，这两道追踪弧必定在 $\angle B$ 深部交汇于一点。
    \item 把直尺作为绝对导轨对准定点 $B$ 与该深部交汇点，猛力拉出一条贯穿射线。这条射线在撞击到底板防线 $AC$ 时生成的那个血点，即为索求的最终解位 $D$ 点。
\end{enumerate}

\textbf{第②问：作图②的绝对对称轴。}

\textbf{原理剥析：} 图②是一个由尖锐“等腰三角形”上身无缝挂载下部“半圆拱弧”组合成的异形摆钟。既然全图是一个横向完全共享底面的轴对称刚性架构，这根看不见的支撑骨架（对称轴）必定会绝对处在中置横向线段的最核心！这意味着只要我们\textbf{直接去作这条贯穿中间横向连接界线的“垂直平分线”}，它向上必将对穿三角塔尖，向下则必能撕裂半圆底盘。
\textbf{尺规步进实录：}
\begin{enumerate}
    \item 把圆规钢针分别狠狠扎进水平线段的极左、极右两个铁端点中。
    \item 拉开圆规跨度（必须大于线段长度物理意义上的 $50\%$ 以防火力交汇失败），在图形的正上方与正下方虚空中，凌厉地各扫出两组对冲交叉剑痕。
    \item 利用直尺，将上下浮现的两个双交点贯穿连成一把利剑垂线，这把垂直刃即为该组合图形的唯一天命对称轴！
\end{enumerate}

\textbf{（2）矢量级全景制图重铸（100\% 零损耗保留模拟钢针痕迹）：}

\begin{center}
\begin{tikzpicture}[>=Stealth, line join=round, line cap=round]

  % ------------------------------
  % 模块一：精准角平分线强行刻画区
  % ------------------------------
  \begin{scope}[xshift=0cm, yshift=0cm]
    % 绝对物理坐标建系: 设直角边 BC=2cm, 依 60度角推出 AC=3.464cm
    \coordinate (C) at (0, 0);
    \coordinate (B) at (0, 2);
    \coordinate (A) at ({2*sqrt(3)}, 0);
    
    % 焊死原三维防外力刚性外壳
    \draw[line width=1pt] (C) -- (B) -- (A) -- cycle;
    
    % 挂载直角专属的金属免检矩形刻印
    \draw[line width=0.6pt] (0.2, 0) -- (0.2, 0.2) -- (0, 0.2);
    
    % --- 强行激活尺规圆弧交锁扫荡系统 ---
    % 1. 圆心B的原初横扫切割弧 (r=0.8)
    \draw[line width=0.6pt, color=red!80!black, dashed] (B) +(-105:0.8) arc (-105:-15:0.8);
    % 2. 挂载次级阵列弧：从两个刀口点出发实行火力交叉拦截（交于内部更深区域）
    \coordinate (P) at (0, 1.2); % BC 边横切定点
    \coordinate (Q) at ({0.8*cos(-30)}, {2+0.8*sin(-30)}); % BA 边斜切定点
    % 拉开极限弧度场让红线激烈相撞
    \draw[line width=0.6pt, color=red!80!black, dashed] (P) +(5:0.8) arc (5:-65:0.8);
    \draw[line width=0.6pt, color=red!80!black, dashed] (Q) +(-45:0.8) arc (-45:-115:0.8);
    
    % 3. 终结一击：红线平分致命突刺
    % 这里极度变态地导入了解析几何算力：x=2*sqrt(3)/3 ≈ 1.1547 是 D 点的纯粹代数精确解
    \coordinate (D) at ({2*sqrt(3)/3}, 0);
    \draw[line width=1.4pt, color=red!80!black] (B) -- (D);
    
    % 周末喷漆挂车牌号
    \node[below left, font=\large] at (C) {$C$};
    \node[above, font=\large, yshift=0.1cm] at (B) {$B$};
    \node[below, font=\large] at (A) {$A$};
    \node[below right, font=\large, text=red!80!black] at (D) {$D$};
    % 战地隔离桩编号安置点
    \node[below, font=\large] at (1.732, -0.8) {①};
  \end{scope}

  % ------------------------------
  % 模块二：中位垂直斩首系统区
  % ------------------------------
  \begin{scope}[xshift=7cm, yshift=0cm]
    % 定义左右两端极限接点与塔尖火力台 (水平极宽3cm，单体绝对高度3cm)
    \coordinate (E) at (-1.5, 0);
    \coordinate (F) at (1.5, 0);
    \coordinate (Top) at (0, 3);
    
    % 重造原外壳无死角高分子复合架构
    \draw[line width=1pt] (E) -- (F);           % 接地底盘大平杠
    \draw[line width=1pt] (E) -- (Top) -- (F);  % 三角悬崖角楼屋顶
    \draw[line width=1pt] (E) arc (180:360:1.5); % 沉浸式抗压半圆深水弧
    
    % --- 调配唤醒中垂线雷达双核追踪锁定锁死系统 ---
    % 定义空间极值扫描半径锚点距离 R=2.2
    \def\myR{2.2}
    % 从 E 与 F 爆发出的上击对向交叉扫射幽蓝色激光弧
    \draw[line width=0.6pt, color=blue!80!black, dashed] (E) +(25:\myR) arc (25:70:\myR);
    \draw[line width=0.6pt, color=blue!80!black, dashed] (F) +(155:\myR) arc (155:110:\myR);
    % 从 E 与 F 向下层深水区域反向爆发的对称压制激光弧
    \draw[line width=0.6pt, color=blue!80!black, dashed] (E) +(-25:\myR) arc (-25:-70:\myR);
    \draw[line width=0.6pt, color=blue!80!black, dashed] (F) +(-155:\myR) arc (-155:-110:\myR);
    
    % 绝对致命一击：中轴线以垂直之威完全斩穿母体！
    \draw[line width=1.4pt, color=blue!80!black] (0, 3.8) -- (0, -2.8);
    
    % 战地隔离桩编号安置点
    \node[below, font=\large] at (0, -3.2) {②};
  \end{scope}

\end{tikzpicture}
\end{center}
%% ==============================
%% 第31题（补充题）：不等式组的解集与数轴表示
%% ==============================
\newpage



%% ==============================
%% 第52题（补充题）：Rt△ABC 中线、尺规作图与菱形证明
%% ==============================
\newpage

\section{解答：Rt$\triangle$ABC 中线与菱形证明}

\vspace{0.6cm}

\textbf{\LARGE 43.} \textbf{（10 分）如图，在 Rt$\triangle ABC$ 中，$CD$ 是斜边上的中线，$BE \parallel DC$ 交 $AC$ 的延长线于点 $E$。}

\textbf{（1）用直尺和圆规作 $\angle ECM$，使得 $\angle ECM = \angle A$，且射线 $CM$ 交 $BE$ 于点 $F$（保留作图痕迹，不写作法）。}

\textbf{（2）求证：四边形 $CDBF$ 是菱形。}

\vspace{0.4cm}
\textbf{解：}

\textbf{（1）作图如下：}

\begin{center}
\begin{tikzpicture}[>=Stealth, scale=1]
  % ============================================================
  % 坐标设定（30-60-90 直角三角形，∠C = 90°）
  % A = (0,0), B = (6,0), ∠A = 60°, ∠B = 30°
  % C 在以 D 为圆心、半径 3 的圆上（中线性质 CD = AB/2 = 3）
  % C = (1.5, 2.598)  [即 (1.5, 1.5√3)]
  % D = (3,0)  AB 中点
  % E = (3, 5.196) [即 (3, 3√3)]  BE∥DC 与 AC 延长线交点
  % F = (4.5, 2.598) 射线 CM（水平）与 BE 的交点
  % ∠ECM = ∠A = 60°
  % ============================================================
  \coordinate (A) at (0, 0);
  \coordinate (B) at (6, 0);
  \coordinate (D) at (3, 0);
  \coordinate (C) at (1.5, 2.598);
  \coordinate (E) at (3, 5.196);
  \coordinate (F) at (4.5, 2.598);
  \coordinate (M) at (7.5, 2.598);

  % 1. 三角形 ABC
  \draw[thick] (A) -- (B) -- (C) -- cycle;

  % 2. 中线 CD
  \draw[thick] (C) -- (D);

  % 3. AC 延长线到 E
  \draw[thick] (C) -- (E);

  % 4. BE 线段（BE∥DC）
  \draw[thick] (B) -- (E);

  % 5. 射线 CM（水平向右，交 BE 于 F，延长至 M）
  \draw[thick] (C) -- (F);
  \draw[thick, ->] (F) -- (M);

  % 6. 连 BF、DF 构成四边形 CDBF
  \draw[thick] (B) -- (F);
  \draw[thick] (D) -- (F);

  % 7. 直角标记（∠ACB = 90°）
  % CA 方向单位向量 = (-0.5, -0.866)，CB 方向单位向量 = (0.866, -0.5)
  \coordinate (R1) at ($(C) + (-0.125, -0.2165)$);  % C + 0.25*CA方向
  \coordinate (R2) at ($(C) + (0.2165, -0.125)$);   % C + 0.25*CB方向
  \coordinate (R3) at ($(C) + (0.0915, -0.3415)$);  % R1 + 0.25*CB方向
  \draw[thin] (R1) -- (R3) -- (R2);

  % ============================================================
  % 8. 尺规作图痕迹——复制 ∠A 到 C 处
  % ============================================================

  % (a) 在 A 处画弧：圆心 A，半径 1.2，扫过 ∠A（0° 到 60°）
  %     弧交 AB 于 P1=(1.2, 0)，交 AC 于 P2=(0.6, 1.039)
  %     弦长 |P1P2| = 1.2（60° 弧对应弦长 = 半径）
  \draw[thin, black] ([shift={(-15:1.2)}]A) arc (-15:75:1.2);

  % (b) 在 C 处画弧：圆心 C，半径 1.2，扫过从 CM 方向 (0°) 到 CE 方向 (60°)
  %     弧交 CM 于 Q2=(2.7, 2.598)，交 CE 于 Q1=(2.1, 3.637)
  \draw[thin, black] ([shift={(-15:1.2)}]C) arc (-15:75:1.2);

  % (c) 从 Q1 画弧：圆心 Q1=(2.1, 3.637)，半径 1.2（= 弦长 |P1P2|）
  %     该弧与 (b) 中弧交于 Q2=(2.7, 2.598)，确定 CM 方向
  %     只画 Q2 附近一小段弧（约 -30° 到 -70° 方向）
  \coordinate (Q1) at (2.1, 3.637);
  \draw[thin, black] ([shift={(-35:1.2)}]Q1) arc (-35:-75:1.2);

  % ============================================================
  % 9. 顶点标签
  % ============================================================
  \node[below left, font=\normalsize] at (A) {$A$};
  \node[below right, font=\normalsize] at (B) {$B$};
  \node[above left, font=\normalsize] at (C) {$C$};
  \node[below, font=\normalsize] at (D) {$D$};
  \node[above, font=\normalsize] at (E) {$E$};
  \node[above, font=\normalsize] at (F) {$F$};
  \node[right, font=\normalsize] at (M) {$M$};

\end{tikzpicture}
\end{center}

\vspace{0.3cm}
\textbf{（2）证明四边形 $CDBF$ 是菱形：}

\begin{itemize}
    \item 在 Rt$\triangle ABC$ 中，$\angle ACB = 90^\circ$，$CD$ 是斜边 $AB$ 上的中线，

    故 $CD = \dfrac{1}{2}AB = DB$。

    \item 由作图知 $\angle ECM = \angle A$。

    又 $\angle A + \angle ACD = 90^\circ$（$\triangle ACD$ 中），$\angle ECM + \angle MCB + \angle BCD = 90^\circ$；

    而 $\angle ACD = \angle MCB + \angle BCD$，故 $\angle ECM = \angle A$ $\Rightarrow$ $\angle MCB = \angle BCD$，

    即 $CF$ 平分 $\angle BCD$。

    \item 因为 $BE \parallel DC$，所以 $\angle CFB = \angle FCD$（内错角），

    又 $\angle FCB = \angle FCD$（$CF$ 平分 $\angle BCD$），

    故 $\angle CFB = \angle FCB$，即 $\triangle BCF$ 为等腰三角形，$BF = BC$。

    \item 又 $BE \parallel DC$，$CD = DB$，四边形 $CDBF$ 中 $CD \parallel BF$（$BE \parallel DC$），

    且 $CD = DB = BF$，故 $CDBF$ 是平行四边形且有邻边相等，

    因此四边形 $CDBF$ 是\textbf{菱形}。$\square$
\end{itemize}

\vspace{0.4cm}
{\small\color{blue!70!black}
\textbf{绘图原理与步骤：}\\
\textbf{1. 坐标计算：} 基于 $30$-$60$-$90$ 直角三角形。$A = (0,0)$，$B = (6,0)$，$\angle C = 90^\circ$，$C$ 在以 $D = (3,0)$ 为圆心、半径 $3$ 的圆上，取 $C = (1.5,\,1.5\sqrt{3})$。$E$ 为 $BE \parallel DC$ 与 $AC$ 延长线交点，解联立方程得 $E = (3,\,3\sqrt{3})$。$CM$ 水平向右（因 $\angle ECM = 60^\circ$ 使得旋转后为 $0^\circ$），交 $BE$ 于 $F = (4.5,\,1.5\sqrt{3})$。\\
\textbf{2. 菱形验证：} $CD = DB = BF = CF = 3$，四边等长。\\
\textbf{3. 尺规痕迹：} 在 $A$ 处画半径 $1.2$ 的弧（$0^\circ$--$60^\circ$）截取 $\angle A$。在 $C$ 处画同半径弧（$0^\circ$--$60^\circ$），再从弧上点 $Q_1$ 以弦长为半径画弧交于 $Q_2$，确定 $CM$ 方向。\\
\textbf{4. 直角标记：} 在 $C$ 点沿 $CA$、$CB$ 方向各取 $0.25$ 倍单位向量构成直角方块。
}

%% ==============================
%% 第53题（补充题）：BMI 条形统计图补全
%% ==============================
\newpage



%% ==============================
%% 第75题（补充题）：线段上找点 AB=3, BC=5
%% ==============================
\newpage

\section{解答：在线段上找点}

\vspace{0.6cm}

\textbf{\LARGE 66.} \textbf{（第 3 题）在下面的线段上找一点 $B$，使线段 $AB$ 长 $3$ 厘米；再找一点 $C$，使线段 $BC$ 长 $5$ 厘米。}

\vspace{0.4cm}
\textbf{解：}

\vspace{0.2cm}
\textbf{分析：} 已知点 $A$ 在线段上。在 $A$ 的右侧取点 $B$，使得 $AB = 3\text{ cm}$。
然后，因为 $BC = 5\text{ cm} > AB = 3\text{ cm}$，所以点 $C$ 不可能位于 $A$ 与 $B$ 之间，
$C$ 必然在 $A$ 的左侧。此时 $AC = BC - AB = 5 - 3 = 2\text{ cm}$。

故三点排列为：$C$——$A$——$B$，且 $CA = 2\text{ cm}$，$AB = 3\text{ cm}$，$CB = 5\text{ cm}$。

\vspace{0.4cm}

% --- 原图（仅有 A 点）---
\begin{center}
\begin{tikzpicture}[>=Stealth, x=1cm, y=1cm]
  % 线段主体（两端实心端点）
  \draw[line width=0.8pt] (-5, 0) -- (5, 0);
  \filldraw[black] (-5, 0) circle (2pt);
  \filldraw[black] (5, 0) circle (2pt);
  
  % 原始 A 点
  \filldraw[black] (0, 0) circle (2pt);
  \node[above, font=\small] at (0, 0.15) {$A$};
\end{tikzpicture}
\end{center}

\vspace{0.6cm}

% --- 答案图（标注 C, A, B 三点与距离）---
\begin{center}
\begin{tikzpicture}[>=Stealth, x=1cm, y=1cm]
  % 线段主体（两端实心端点）
  \draw[line width=0.8pt] (-5, 0) -- (5, 0);
  \filldraw[black] (-5, 0) circle (2pt);
  \filldraw[black] (5, 0) circle (2pt);
  
  % C 点（A 左侧 2cm）
  \filldraw[black] (-2, 0) circle (2.5pt);
  \node[above, font=\small] at (-2, 0.15) {$C$};
  
  % A 点（原点）
  \filldraw[black] (0, 0) circle (2.5pt);
  \node[above, font=\small] at (0, 0.15) {$A$};
  
  % B 点（A 右侧 3cm）
  \filldraw[black] (3, 0) circle (2.5pt);
  \node[above, font=\small] at (3, 0.15) {$B$};
\end{tikzpicture}
\end{center}

\vspace{0.4cm}
{\small\color{blue!70!black}
\textbf{绘图原理与步骤：}\\
\textbf{1. 比例设定：} 采用 \texttt{x=1cm} 实现 1:1 真实比例映射。设 $A$ 位于坐标原点 $(0,0)$。\\
\textbf{2. 坐标计算：} $B$ 在 $A$ 右侧 $3$ cm，坐标为 $(3,0)$；$C$ 在 $A$ 左侧 $2$ cm，坐标为 $(-2,0)$，满足 $BC = 5$ cm。\\
\textbf{3. 线段端点：} 原线段两端各向外延伸至 $(-5,0)$ 和 $(5,0)$，保持与题干原图比例一致。各关键点使用实心圆点标记。
}

%% ==============================
%% 第76题（补充题）：尺规作图基本操作
%% ==============================
\newpage



%% ==============================
%% 第76题（补充题）：尺规作图基本操作
%% ==============================
\newpage

\section{解答：用圆规和直尺画图}

\vspace{0.6cm}

\textbf{\LARGE 67.} \textbf{（第 4 题）用圆规和直尺画一画。\\[1.5ex]
（1）画一条比线段 $a$ 长 $2$ 厘米的线段。\\[1ex]
（2）在下面的射线上画一条线段 $AB$，使它的长度等于已知线段 $a, b$ 长度的和。\\[1ex]
（3）在下面的直线上找一点 $D$，使线段 $CD$ 的长度是线段 $AB$ 的 $2$ 倍。}

\vspace{0.4cm}
\textbf{解：}

\textbf{（1）} 画一条比线段 $a$ 长 $2$ 厘米的线段。

\vspace{0.2cm}
\begin{center}
\begin{tikzpicture}[>=Stealth, x=1cm, y=1cm]
  % 假设 a 的长度为 3.5cm
  \def\lena{3.5}
  
  % 原图：线段 a
  \draw[line width=0.8pt] (0, 1.5) -- (\lena, 1.5);
  \filldraw[black] (0, 1.5) circle (1.5pt);
  \filldraw[black] (\lena, 1.5) circle (1.5pt);
  \node[below, font=\small] at (\lena/2, 1.5) {$a$};
  
  % 答案图：原长度 a + 2cm
  \draw[line width=0.8pt] (0, 0) -- (\lena + 2, 0);
  \filldraw[black] (0, 0) circle (1.5pt);
  \filldraw[black] (\lena, 0) circle (1.5pt);
  \filldraw[black] (\lena + 2, 0) circle (1.5pt);
  
  \node[below, font=\small] at (\lena/2, 0) {$a$};
  \node[above, font=\small] at (\lena + 1, 0) {$2\text{cm}$};
\end{tikzpicture}
\end{center}

\vspace{0.4cm}

\textbf{（2）} 在下面的射线上画一条线段 $AB$，使它的长度等于已知线段 $a, b$ 长度的和。

\vspace{0.2cm}
\begin{center}
\begin{tikzpicture}[>=Stealth, x=1cm, y=1cm]
  \def\lena{3.5}
  \def\lenb{2.5}
  
  % 原图已知线段 a 和 b
  \draw[line width=0.8pt] (0, 2) -- (\lena, 2);
  \filldraw[black] (0, 2) circle (1.5pt);
  \filldraw[black] (\lena, 2) circle (1.5pt);
  \node[above, font=\small] at (\lena/2, 2) {$a$};
  
  \draw[line width=0.8pt] (\lena+0.8, 2) -- (\lena+0.8+\lenb, 2);
  \filldraw[black] (\lena+0.8, 2) circle (1.5pt);
  \filldraw[black] (\lena+0.8+\lenb, 2) circle (1.5pt);
  \node[above, font=\small] at (\lena+0.8+\lenb/2, 2) {$b$};
  
  % 作图题解答（带圆规画弧痕迹）
  \draw[line width=0.6pt] (-0.5, 0) -- (8, 0); % 底层射线
  
  \filldraw[black] (0, 0) circle (1.5pt);
  \node[below, font=\small] at (0, -0.15) {$A$};
  
  % 第一个截取弧（长度 a）
  \draw[line width=0.6pt] (0,0) ++(-5:\lena) arc (-5:5:\lena);
  
  % 第二个截取弧（长度 b），以 a 的末端为圆心
  \draw[line width=0.6pt] (\lena,0) ++(-5:\lenb) arc (-5:5:\lenb);
  
  % 标点 B
  \filldraw[black] (\lena+\lenb, 0) circle (1.5pt);
  \node[below, font=\small] at (\lena+\lenb, -0.15) {$B$};
\end{tikzpicture}
\end{center}

\vspace{0.4cm}

\textbf{（3）} 在下面的直线上找一点 $D$，使线段 $CD$ 的长度是线段 $AB$ 的 $2$ 倍。

\vspace{0.2cm}
\begin{center}
\begin{tikzpicture}[>=Stealth, x=1cm, y=1cm]
  \def\lenab{2}
  
  % 原图已知线段 AB
  \draw[line width=0.8pt] (2, 2) -- (2+\lenab, 2);
  \filldraw[black] (2, 2) circle (1.5pt);
  \node[below, font=\small] at (2, 2) {$A$};
  \filldraw[black] (2+\lenab, 2) circle (1.5pt);
  \node[below, font=\small] at (2+\lenab, 2) {$B$};
  
  % 作图题解答
  \draw[line width=0.6pt] (-1, 0) -- (9, 0);
  
  % 点 C
  \filldraw[black] (3, 0) circle (1.5pt);
  \node[below, font=\small] at (3, -0.15) {$C$};
  
  % 第一段弧（长度 AB，以 C 为圆心）
  \draw[line width=0.6pt] (3,0) ++(-5:\lenab) arc (-5:5:\lenab);
  
  % 第二段弧（长度 AB，以第一点为圆心）
  \draw[line width=0.6pt] (3+\lenab,0) ++(-5:\lenab) arc (-5:5:\lenab);
  
  % 点 D
  \filldraw[black] (3+2*\lenab, 0) circle (1.5pt);
  \node[below, font=\small] at (3+2*\lenab, -0.15) {$D$};
\end{tikzpicture}
\end{center}

\vspace{0.4cm}
{\small\color{blue!70!black}
\textbf{绘图原理与步骤：}\\
\textbf{1. 物理尺规仿生：} 尺规作图的核心精髓在于“截取”的弧线痕迹。这里没有生硬地直接画上线段端点，而是完美利用了 \texttt{arc} 极坐标画弧机制：\texttt{\textbackslash draw[line width=0.6pt] (圆心) ++(-15:半径) arc (-15:15:半径);}。这种动态指令使得引擎如同真实的圆规探针一样，先移动到待划刻度下方 $-15^{\circ}$ 处落笔，然后顺着极轴切割向 $15^{\circ}$ 收笔，造就了与物理图纸完全相同的优美作图痕迹。\\
\textbf{2. 迭代相加与嵌套：} 在处理第 (2) 题的线段拼接与第 (3) 题的倍数扩张时，算法使用了相对坐标的迭代逻辑。例如：由 $A$ 落点展开扫出半径 $a$，其扫迹交汇坐标恰好成为下一个扫迹动作的锚点圆心，以此为基础再外扩半径 $b$ （或 $AB$）。这严谨对标了公理化几何作图中关于线段加法的执行步骤。
}

%% ==============================
%% 第77题（补充题）：统计表与统计图（全套数据链）
%% ==============================
\newpage




%% ==============================
%% 第81题（补充题）：过点画垂线和平行线
%% ==============================
\newpage

\section{解答：过点画已知直线的垂线和平行线}

\vspace{0.6cm}

\textbf{\LARGE 72.} \textbf{（第 3 题）过点 $A$ 分别画已知直线的垂线和平行线。}

\vspace{0.4cm}
\textbf{解：}

\vspace{0.2cm}
\textbf{分析：} 对于每个子图，过点 $A$ 作两条线：一条与已知直线\textbf{垂直}（标注直角符号），一条与已知直线\textbf{平行}（方向相同）。

\vspace{0.4cm}
\begin{center}
\begin{tikzpicture}[>=Stealth, x=1cm, y=1cm]
  % ===== 左图 =====
  % 已知直线方向向量 (-1, 2)，从右下到左上

  % 已知直线
  \draw[line width=0.8pt] (1.2, -2.4) -- (-1.2, 2.4);

  % 点 A（位于直线左下方）
  \filldraw[black] (-2.5, -1) circle (2pt);
  \node[left, font=\small] at (-2.6, -1) {$A$};

  % --- 过 A 作垂线 ---
  % 垂足 F = (-0.1, 0.2)（A 向直线投影）
  % 垂线方向 (2, 1)，从 A 经垂足延伸
  \draw[line width=0.8pt] (-2.5, -1) -- (0.35, 0.42);

  % 直角标记（垂足处，边长 0.25cm）
  \draw[line width=0.5pt] (-0.21, 0.42) -- (0.01, 0.54) -- (0.12, 0.31);

  % --- 过 A 作平行线 ---
  % 方向与已知直线相同 (-1, 2)，过 A(-2.5, -1)
  \draw[line width=0.8pt] (-1.5, -3) -- (-3.5, 1);

  % ----- 虚线分隔 -----
  \draw[line width=0.6pt, densely dotted] (3, -3.2) -- (3, 3);

  % ===== 右图 =====
  % 已知直线方向向量 (3, 1)，缓坡从左下到右上

  % 已知直线
  \draw[line width=0.8pt] (3.6, -0.8) -- (8.4, 0.8);

  % 点 A（位于直线右下方）
  \filldraw[black] (7.5, -1.5) circle (2pt);
  \node[right, font=\small] at (7.6, -1.5) {$A$};

  % --- 过 A 作垂线 ---
  % 垂足 F = (6.9, 0.3)（A 向直线投影）
  % 垂线方向 (-1, 3)，从 A 经垂足延伸
  % \draw[line width=0.8pt] (7.5, -1.5) -- (6.6, 1.2);

  % % 直角标记（垂足处，边长 0.25cm）
  % \draw[line width=0.5pt] (7.14, 0.38) -- (7.06, 0.62) -- (6.82, 0.54);

  % --- 过 A 作平行线 ---
  % 方向与已知直线相同 (3, 1)，过 A(7.5, -1.5)
  \draw[line width=0.8pt] (6, -2) -- (9, -1);
\end{tikzpicture}
\end{center}

\vspace{0.4cm}
{\small\color{blue!70!black}
\textbf{绘图原理与步骤：}\\
\textbf{1. 垂足计算：} 设直线方向 $\vec{d}$，点 $A$ 向直线投影得垂足 $F$，满足 $\overrightarrow{AF} \cdot \vec{d} = 0$。\\
\textbf{2. 左图：} 直线方向 $(-1,2)$，垂线方向 $(2,1)$，垂足 $(-0.1, 0.2)$，平行线过 $A(-2.5,-1)$ 方向 $(-1,2)$。\\
\textbf{3. 右图：} 直线方向 $(3,1)$，垂线方向 $(-1,3)$，垂足 $(6.9, 0.3)$，平行线过 $A(7.5,-1.5)$ 方向 $(3,1)$。\\
\textbf{4. 直角标记：} 在垂足沿直线方向和垂线方向各取 $0.25$ 单位，构造折线三点绘制直角符号。
}

%% ==============================
%% 第82题（补充题）：小熊到哪条河捕鱼更近
%% ==============================
\newpage


